% from ega4-2-fr.tex
\oldpage[IV-2]{27}
tout $i$,
$(B_i\otimes_{k_i}k'_i)_s$ et $W_i\otimes_{k_i}k'_i$ sont des
$(B_i\otimes_{k_i}k'_i)$-modules isomorphes~; il suffit donc de prouver
(2.5.8) lorsque $n=1$, c'est-à-dire dans le cas où $A$ et $A'$ sont des
corps, $B$ une algèbre simple de centre $A$.
\end{env}

\begin{env}[2.5.8.9]
\label{IV.2.5.8.9-fr}
\emph{Fin de la démonstration.} On sait alors (Bourbaki, \emph{Alg.},
chap.~VIII, \S~5, prop.~6 et 8) que tout $B$-module est somme directe de
modules isomorphes à un idéal minimal de $B$, et deux $B$-modules de rang
fini sur $A$ sont donc isomorphes si et seulement s'ils ont même rang sur
$A$. Par hypothèse, on a $[W':A']=[B'_s:A']$. Mais on a
$[W':A']=[W:A]$ et $[B'_s:A']=[B_s:A]$~; donc $[W:A]=[B_s:A]$, ce qui
termine la démonstration.
\end{env}

\subsection{Permanence de propriétés ensemblistes et topologiques de morphismes
par descente fidèlement plate}
\label{subsection:IV.2.6-fr}

\begin{proposition}[2.6.1]
\EGAReaderTarget{ega4.c04.prop.2-6-1}{IV.2.6.1-fr}
Soient $f:X\to Y$ un $S$-morphisme de $S$-préschémas, $g:S'\to S$ un
morphisme surjectif, $X'=X_{(S')}$, $Y'=Y_{(S')}$,
$f'=f_{(S')}:X'\to Y'$. Considérons, pour un morphisme, la propriété d'être~:
\begin{enumerate}
\item[(i)] surjectif~;
\item[(ii)] injectif~;
\item[(iii)] à fibres finies (en tant qu'ensembles)~;
\item[(iv)] bijectif~;
\item[(v)] radiciel.
\end{enumerate}
Alors, si $P$ désigne une des propriétés précédentes et si $f'$ possède la
propriété $P$, il en est de même de $f$.
\end{proposition}

Comme le morphisme projection $Y'\to Y$ est lui aussi surjectif (I, 3.5.2),
on peut, en vertu de (I, 3.3.11), se borner au cas où $Y=S$, $Y'=S'$. Pour
tout $y\in Y$ (resp. $y'\in Y'$) désignons par $X_y$ (resp. $X'_{y'}$) le
préschéma fibre $f^{-1}(y)$ (resp. $f'^{-1}(y')$) (I, 3.6.2)~; on sait que
pour $y'\in Y'$, $y=g(y')$, on a un isomorphisme canonique
\[
X'_{y'}\xrightarrow{\sim}X_y\otimes_{k(y)}k(y')\quad\text{(I, 3.6.4)}~;
\]
comme le morphisme $\operatorname{Spec}(k(y'))\to\operatorname{Spec}(k(y))$
est surjectif, il en est de même de la projection $X'_{y'}\to X_y$
(I, 3.5.2)~; donc si $X'_{y'}$ est non vide (resp. a au plus un point, resp.
est un ensemble fini), il en est de même de $X_y$~; comme $g:Y'\to Y$ est
surjectif, cela démontre la proposition dans les cas (i), (ii) et (iii), et
(iv) résulte de (i) et (ii)~; enfin, pour prouver (v), il suffit de montrer
que si $f'$ est universellement injectif (I, 3.5.11), il en est de même de
$f$~; or, soit $Y_1\to Y$ un morphisme quelconque~; posons
$X_1=X\times_Y Y_1$ et $f_1=f_{(Y_1)}$. D'autre part, posons
\[
Y'_1=Y_1\times_Y Y',\quad X'_1=X'\times_{Y'}Y'_1=X_1\times_{Y_1}Y'_1,
\quad f'_1=f'_{(Y'_1)}=(f_1)_{(Y'_1)}:X'_1\to Y'_1~;
\]
comme $g'=g_{(Y_1)}:Y'_1\to Y_1$ est surjectif (I, 3.5.2) et que $f'$ est
universellement injectif, $f'_1$ est injectif, et il résulte de (ii) que
$f_1$ est injectif, d'où notre assertion.

\begin{proposition}[2.6.2]
\label{IV.2.6.2-fr}
Les notations étant celles de (2.6.1), supposons le morphisme $g:S'\to S$
fidèlement plat et quasi-compact. Considérons, pour un morphisme, la propriété
d'être~:
\begin{enumerate}
\item[(i)] ouvert~;
\item[(ii)] fermé~;
\item[(iii)] quasi-compact et un homéomorphisme sur le sous-espace image~;
\item[(iv)] un homéomorphisme.
\end{enumerate}
Alors, si $P$ désigne une des propriétés précédentes et si $f'$ possède la
propriété $P$, il en est de même de $f$.
\end{proposition}

